\chapter{Architecture et réalisation}
\label{chap:archi}

Ce chapitre constitue le c{\oe}ur du travail d'ingénierie. Il détaille
l'architecture globale et par couche, les pipelines RAG, RAGOps, de traçabilité
et de l'AI Execution Studio, les flux d'exécution, ainsi que les aspects
transverses de sécurité, de performance, de scalabilité et d'observabilité.

\section{Architecture globale}
\devpilot{} est un système orienté services déployé par conteneurs. La vision
architecturale repose sur quatre principes : séparation des responsabilités en
couches, asynchronisme pour découpler le chemin critique du traitement lourd,
traçabilité native (chaque étape produit une trace persistée) et neutralité
vis-à-vis des fournisseurs d'IA.

\paragraph{Explication.} Le diagramme suivant présente le flux de bout en bout à
travers les composants principaux.
\paragraph{Hypothèses.} Les magasins de données sont distingués ; le pipeline
agentique est synthétisé en un composant.

\begin{lstlisting}[style=eraser,caption={Architecture globale et flux de données (Eraser)},label={lst:global}]
colorMode bold
styleMode shadow
direction down

User [icon: user]
Client [icon: monitor, color: blue] { Frontend [icon: nextjs] }
API Gateway [icon: server, color: blue] { FastAPI [icon: python]; Auth [icon: lock] }
Orchestration [icon: workflow, color: purple] { Workflow [icon: cpu, label: "7 agents"] }
Async [icon: refresh-cw, color: orange] { Redis [icon: redis]; Celery [icon: python] }
Data [icon: database, color: green] { Postgres [icon: postgresql]; Qdrant [icon: database] }
AI [icon: cpu, color: teal] { Ollama [icon: server]; Gemini [icon: google] }
Obs [icon: activity, color: red] { Prometheus [icon: prometheus]; Grafana [icon: grafana] }

User > Frontend: question
Frontend > FastAPI: HTTPS / SSE
FastAPI > Workflow
FastAPI > Redis: enqueue
Redis > Celery
Workflow > Qdrant: recherche
Workflow > Ollama: embeddings
Workflow > Gemini: génération
Celery > Qdrant: upsert
FastAPI > Postgres
Celery > Postgres
FastAPI > Prometheus
Prometheus > Grafana
\end{lstlisting}

\paragraph{Interprétation.} La requête synchrone traverse l'API et le pipeline ;
l'ingestion documentaire est déléguée à Celery via Redis. PostgreSQL conserve
l'état et les traces, Qdrant les vecteurs.
\eraserfig{Architecture système globale}{fig:global}

\section{Architecture Backend}
Le \emph{backend} FastAPI s'organise en couches aux frontières nettes :
\texttt{api} (routes et intégration HTTP), \texttt{core} (configuration, sécurité,
exceptions), \texttt{db} (session et base ORM), \texttt{models} (modèles
SQLAlchemy), \texttt{schemas} (schémas Pydantic), \texttt{services} (logique
applicative), \texttt{agents} (orchestration agentique), \texttt{workers} (tâches
Celery) et \texttt{utils}. Les routes restent minces et délèguent aux services, ce
qui isole la logique métier et facilite les tests.

\section{Architecture Frontend}
Le \emph{frontend} Next.js (\emph{App Router}, React~19) adopte une organisation
par modules fonctionnels (\texttt{features/}), un système de design fondé sur des
jetons CSS (« Horizon », thème clair/sombre) et la gestion de l'état serveur par
React Query. Les visualisations lourdes (AI Execution Studio, Architecture
Explorer) sont chargées de manière différée (\emph{lazy loading}) afin de
préserver le poids initial des routes.

\section{Architecture Base de données}
Le modèle relationnel comprend onze tables gérées par SQLAlchemy et versionnées
par Alembic. \devpilot{} repose sur un \textbf{stockage dual} : PostgreSQL
conserve le texte canonique des fragments et une référence \texttt{vector\_id}
vers le point Qdrant correspondant, tandis que Qdrant stocke les vecteurs
indexés en HNSW (distance cosinus). Cette séparation préserve l'intégrité
référentielle tout en confiant la recherche vectorielle à un moteur spécialisé.

\section{Architecture API}
L'API est versionnée sous \texttt{/api/v1} et organisée en routeurs
(authentification, espaces, documents, chat, système, administration, RAGOps,
traces). Elle expose une interface REST pour les opérations CRUD et une interface
\emph{streaming} (SSE) pour la génération. L'autorisation repose sur des
dépendances FastAPI appliquant le contrôle d'accès par rôles.

\section{Architecture RAG}
La génération suit un pipeline en plusieurs étapes : \textbf{ingestion} (extraction
du texte), \textbf{découpage} en fragments à fenêtres chevauchantes,
\textbf{vectorisation} (\emph{embeddings} \texttt{nomic-embed-text}, 768
dimensions), \textbf{indexation} dans Qdrant, \textbf{récupération} hybride
(sémantique et lexicale) d'un ensemble large de candidats, \textbf{reclassement}
fin, \textbf{construction du prompt} (système + contexte + question) et
\textbf{génération} ancrée et citée. La récupération suit une stratégie « rappel
large puis précision fine » : quinze candidats sont remontés puis réduits aux cinq
meilleurs.

\paragraph{Explication.} Le diagramme suivant détaille le flux agentique de sept
agents coopérants.
\paragraph{Hypothèses.} Chaque agent est isolé et persisté ; le correcteur
n'intervient qu'en cas de réponse non fidèle.

\begin{lstlisting}[style=eraser,caption={Pipeline agentique RAG (Eraser)},label={lst:pipe}]
colorMode bold
styleMode shadow
direction right

Question [icon: message-circle, color: blue]
Pipeline [icon: workflow, color: purple] {
  Router [icon: git-branch]; QueryRewriter [icon: edit]; Retrieval [icon: search]
  Reranker [icon: bar-chart]; Generator [icon: sparkles]; Evaluator [icon: shield]; Corrector [icon: tool]
}
Qdrant [icon: database, color: green, label: "k=15"]
Gemini [icon: google, color: teal]
Answer [icon: check-circle, color: green, label: "ancrée + citée"]

Question > Router > QueryRewriter > Retrieval
Retrieval > Qdrant: 15 candidats
Retrieval > Reranker > Generator
Generator > Gemini: stream
Generator > Evaluator
Evaluator > Corrector: si non fidèle
Evaluator > Answer: si fidèle
Corrector > Answer
\end{lstlisting}

\paragraph{Interprétation.} Le flux est linéaire jusqu'à l'évaluation ; le
correcteur constitue une boucle de réparation conditionnelle garantissant la
fidélité de la réponse diffusée.
\eraserfig{Pipeline de génération agentique}{fig:pipe}

\section{Architecture RAGOps}
RAGOps constitue la couche d'\emph{exploitation} du pipeline RAG — le pont entre
architecture et observabilité. Là où un tableau de bord classique rapporte des
métriques métier, RAGOps rapporte la santé et la qualité du système d'IA lui-même.
Il agrège plusieurs dimensions :

\begin{description}[leftmargin=2.6cm,style=nextline]
\item[Health Monitoring] Santé de la collection Qdrant et état des services
(\texttt{/admin/ragops/qdrant/health}, \texttt{ops-health}, \texttt{vector-health}).
\item[Retrieval Analytics] Stratégie de récupération, nombre de candidats et
qualité du contexte récupéré.
\item[Evaluation Intelligence] Fidélité moyenne, comptage des réponses à risque
d'hallucination et des réponses corrigées.
\item[Embedding Analytics] Couverture d'\emph{embedding} par espace
(\texttt{chunks\_with\_vector\_id} / \texttt{total\_chunks}).
\item[Workspace Health] Score de santé RAG par espace (0--100) combinant ratio
d'indexation, couverture, fidélité et fiabilité.
\item[Agent Monitoring] Latence moyenne par agent et comptabilité des jetons et
coûts.
\item[Failure Center] Documents en échec et relance de l'ingestion
(\texttt{POST /admin/ragops/documents/\{id\}/retry}).
\item[Observability] Métriques Prometheus et tableaux de bord Grafana.
\end{description}

RAGOps est essentiel dans les systèmes d'IA d'entreprise car la qualité d'un
pipeline RAG dérive silencieusement (documents non indexés, couverture partielle,
hausse des hallucinations) : sans observabilité spécifique à l'IA, ces régressions
restent invisibles jusqu'à l'incident.

\section{Architecture Traceability}
La traçabilité est modélisée comme une chaîne dirigée centrée sur le message
assistant.

\paragraph{Explication.} Le diagramme présente la lignée d'exécution persistée.
\paragraph{Hypothèses.} Chaque message assistant agrège ses exécutions d'agents,
ses passages, son évaluation et sa comptabilité d'usage.

\begin{lstlisting}[style=eraser,caption={Chaîne de traçabilité (Eraser)},label={lst:trace}]
colorMode bold
styleMode shadow
direction right

Message [icon: message-square, color: blue]
Lineage [icon: layers, color: purple] {
  AgentRuns [icon: cpu, label: "agent_runs"]
  RetrievedChunks [icon: file-text, label: "retrieved_chunks"]
  Evaluations [icon: shield, label: "evaluations"]
  LLMUsage [icon: dollar-sign, label: "llm_usage"]
}
Storage [icon: database, color: green, label: "PostgreSQL"]
Inspector [icon: search, color: teal, label: "Trace Inspector"]

Message > AgentRuns: 1..*
Message > RetrievedChunks: 1..*
Message > Evaluations: 1..1
Message > LLMUsage: 0..*
AgentRuns > Storage
Evaluations > Storage
Storage > Inspector: reconstruction
\end{lstlisting}

\paragraph{Interprétation.} Toute réponse est \emph{reconstructible a posteriori}
à partir de ces données, sans dépendance à des journaux externes. Cela fonde
quatre propriétés : \textbf{auditabilité} (justifier une réponse), \textbf{reproductibilité}
(rejouer et comparer), \textbf{explicabilité} (exposer le raisonnement) et
\textbf{débogage/gouvernance} (diagnostiquer les régressions et tracer les actions
d'administration via le journal d'audit).
\eraserfig{Chaîne de traçabilité}{fig:trace}

\section{Architecture AI Execution Studio}
L'AI Execution Studio transforme le pipeline RAG en une expérience animée en temps
réel : au lieu d'afficher des métriques, il donne à voir le raisonnement du système.
Il est piloté par les \emph{mêmes API réelles} que le chat (flux SSE + endpoints de
trace). Il visualise séquentiellement les dix étapes : question, \emph{embedding},
recherche vectorielle, récupération des fragments, reclassement, construction du
contexte, construction du prompt, génération (\emph{streaming} de jetons),
évaluation (jauges animées) et persistance de la trace. Sa valeur est double :
\textbf{pédagogique} (comprendre le RAG en une trentaine de secondes) et
\textbf{d'observabilité} (visualiser où se dépensent la latence et la qualité).

\section{Flux d'ingestion}
L'ingestion est entièrement asynchrone (cf. figure~\ref{fig:seq-ingest}) : l'API
persiste le document avec le statut \texttt{pending} et délègue à Celery la tâche
\texttt{process\_document\_task}, relancée jusqu'à trois fois en cas d'échec, avant
passage au statut \texttt{failed}.

\section{Flux de recherche}
La recherche vectorise la requête réécrite, interroge Qdrant pour un ensemble de
candidats, fusionne signaux sémantiques et lexicaux (récupération hybride), puis
reclasse pour ne conserver que les passages les plus pertinents dans la limite du
budget de contexte (6000 caractères).

\section{Flux de génération}
Le générateur construit le prompt final (système + contexte + question) et produit
une réponse ancrée. Les fragments retenus sont reliés à des marqueurs de citation
\texttt{[n]} ; un mécanisme garantit la présence d'au moins une citation lorsque
du contexte est disponible.

\section{Flux de streaming SSE}
La diffusion repose sur \emph{Server-Sent Events} : le \emph{backend} renvoie une
\texttt{StreamingResponse} (\texttt{text/event-stream}, en-tête
\texttt{X-Accel-Buffering: no}) émettant une séquence d'événements typés
(\texttt{start}, \texttt{trace}, \texttt{token}, \texttt{citations},
\texttt{evaluation}, \texttt{correction}, \texttt{message}, \texttt{done},
\texttt{error}). Le \emph{streaming} repose sur l'API \texttt{streamGenerateContent}
de Gemini ; côté client, le flux est consommé via \texttt{fetch} et un
\texttt{ReadableStream}. Le listing~\ref{lst:sse} illustre le générateur côté
serveur.

\begin{lstlisting}[style=code,language=Python,caption={Génération du flux SSE (extrait)},label={lst:sse}]
async def event_stream() -> AsyncIterator[str]:
    async for event, data in chat_service.stream_chat_events(...):
        yield sse_event(event, data)

return StreamingResponse(
    event_stream(),
    media_type="text/event-stream",
    headers={"X-Accel-Buffering": "no"},
)
\end{lstlisting}

\section{Sécurité}
La sécurité est multicouche : authentification JWT, autorisation par rôles,
isolation des données par espace de travail (toute requête est cantonnée aux
espaces dont l'utilisateur est membre), journalisation des actions
d'administration et masquage systématique des secrets (clés contenant
\emph{password}, \emph{secret}, \emph{token}, \emph{api\_key}) dans les vues de
débogage et de traçabilité.

\section{Authentification JWT}
L'authentification repose sur des jetons JWT signés, émis à la connexion et
transmis dans l'en-tête \texttt{Authorization: Bearer}. Les dépendances FastAPI
résolvent l'utilisateur courant et vérifient son activité avant tout accès aux
ressources protégées.

\section{Gestion des permissions}
Trois rôles structurent l'accès : \texttt{user}, \texttt{admin} et
\texttt{super\_admin}. Le contrôle s'applique à deux niveaux : par route (les
\emph{endpoints} d'administration, RAGOps et traces exigent un rôle privilégié) et
par espace de travail (l'appartenance conditionne l'accès aux documents et
conversations).

\section{Performance}
Les leviers de performance sont l'asynchronisme bout-en-bout (E/S non bloquantes,
déport du calcul lourd vers Celery), la recherche HNSW sous-linéaire, le
\emph{streaming} (réduction du \emph{time-to-first-token}), ainsi que le
chargement différé et le découpage des \emph{bundles} côté \emph{frontend}
(routes de 110 à 175~kB environ).

\section{Scalabilité}
Le \emph{backend} sans état est horizontalement extensible derrière un répartiteur
de charge ; les \emph{workers} Celery se multiplient pour absorber les pics
d'ingestion ; Qdrant supporte le partitionnement et PostgreSQL la réplication en
lecture.

\section{Observabilité}
L'observabilité opère à trois niveaux complémentaires : métriques système
(Prometheus, tableaux Grafana), exploitation applicative (RAGOps) et traçabilité
des requêtes (explorateur et inspecteur de traces), auxquels s'ajoute la
visualisation pédagogique de l'AI Execution Studio.

\section{Conclusion}
L'architecture de \devpilot{} traduit ses principes directeurs en un système
concret, en couches, asynchrone, traçable et neutre vis-à-vis des fournisseurs
d'IA. Les pipelines documentaire, de récupération, de génération agentique et
d'évaluation, instrumentés par RAGOps et la traçabilité, forment un tout cohérent.
Le chapitre suivant justifie les choix technologiques et détaille l'implémentation.
